% *** ก่อนส่ง: ลบข้อความตัวอย่าง (\ldots) ทั้งหมด แล้วเขียนเนื้อหาของตนเอง ***
% =============================================================
% บทที่ 1 - บทนำ
% แก้ไขเนื้อหาในบทนี้ได้ที่ไฟล์ chapter1.tex
% =============================================================

% -----------------------------------------------------------
\section{ความเป็นมาและความสำคัญของปัญหา}
% -----------------------------------------------------------

ในยุคปัจจุบัน ธุรกิจการเช่าอ่านหนังสือการ์ตูนทั้งในรูปแบบร้านค้าและแพลตฟอร์มดิจิทัลต่างมุ่งเน้นการใช้เทคโนโลยีเพื่อสร้างประสบการณ์ที่ดีให้กับผู้ใช้งาน ระบบแนะนำ (Recommender Systems) จึงกลายเป็นเครื่องมือสำคัญในการจัดการกับข้อมูลจำนวนมหาศาล และช่วยให้ผู้ใช้สามารถค้นพบเนื้อหาที่สอดคล้องกับความสนใจของตนได้อย่างมีประสิทธิภาพ \cite{ricci2015}

อย่างไรก็ตาม ระบบแนะนำที่ใช้กันโดยทั่วไปมักอาศัยการวิเคราะห์จาก พฤติกรรมเพียงประเภทเดียว (Single-Behavior) โดยส่วนใหญ่มุ่งเน้นไปที่พฤติกรรมสุดท้ายของกระบวนการตัดสินใจ เช่น การเช่า หรือการซื้อ ซึ่งเรียกว่า Conversion Behavior การพิจารณาข้อมูลเพียงมิติเดียวนี้ก่อให้เกิดปัญหาสำคัญหลายประการ ได้แก่ ปัญหาข้อมูลเบาบาง (Data Sparsity) เนื่องจากผู้ใช้ส่วนใหญ่มีจำนวนการเช่าหรือซื้อไม่มากพอ ไม่สามารถสะท้อนความสนใจที่แท้จริงของผู้ใช้ได้ครบถ้วน เพราะละเลยพฤติกรรมในขั้นตอนก่อนหน้า และความแม่นยำของระบบแนะนำลดลง โดยเฉพาะกับผู้ใช้ใหม่หรือผู้ใช้ที่ยังไม่เคยมีประวัติการเช่า (Cold-Start Problem) \cite{ricci2015}

ในความเป็นจริง ก่อนที่ผู้ใช้จะตัดสินใจเช่าหนังสือการ์ตูน มักมีปฏิสัมพันธ์หลายรูปแบบ เช่น การกดดูรายละเอียด, การค้นหา, การบันทึกไว้ในรายการโปรด, การให้คะแนน และการเขียนรีวิว พฤติกรรมเหล่านี้แม้จะไม่ใช่การเช่าจริง แต่ถือเป็น สัญญาณแฝงของความสนใจ (Latent Intentions) ที่สะท้อนแนวโน้มในการตัดสินใจในอนาคตได้อย่างมีนัยสำคัญ \cite{wang2019dmran}

ด้วยเหตุนี้ แนวคิดของระบบแนะนำแบบหลายพฤติกรรม (Multi-Behavior Recommender System) จึงถูกนำเสนอเพื่อแก้ไขข้อจำกัดของระบบแบบเดิม โดยคำว่า Multi-Behavior หมายถึง การออกแบบระบบแนะนำที่นำข้อมูลจากพฤติกรรมหลายประเภทของผู้ใช้มาวิเคราะห์ร่วมกัน ไม่ว่าจะเป็นพฤติกรรมหลัก (Target Behavior) เช่น การเช่า และพฤติกรรมเสริม (Auxiliary Behaviors) เช่น การดู การค้นหา หรือการบันทึก เพื่อสร้างแบบจำลองความสนใจของผู้ใช้ได้อย่างครอบคลุมมากขึ้น \cite{wang2019dmran}

การใช้แนวคิด Multi-Behavior ช่วยให้ระบบสามารถ ลดปัญหาข้อมูลเบาบาง ด้วยการเพิ่มแหล่งข้อมูลปฏิสัมพันธ์ เข้าใจลำดับขั้นของการตัดสินใจของผู้ใช้ ได้ดียิ่งขึ้น อีกทั้งยังเพิ่มความแม่นยำและความเกี่ยวข้องของคำแนะนำ และรองรับผู้ใช้ใหม่ได้ดีกว่าระบบที่พึ่งพาเฉพาะข้อมูลการเช่า \cite{mbgcn_nd}

ดังนั้น โครงงานนี้จึงมุ่งพัฒนาแอปพลิเคชันระบบแนะนำหนังสือการ์ตูนโดยประยุกต์ใช้แนวคิด Multi-Behavior Recommender System ร่วมกับเทคนิคด้าน Machine Learning เพื่อผสานข้อมูลปฏิสัมพันธ์ที่หลากหลายของผู้ใช้เข้าด้วยกัน พร้อมกำหนดน้ำหนักความสำคัญของแต่ละพฤติกรรมอย่างเหมาะสม อันจะนำไปสู่การสร้างระบบแนะนำที่มีความแม่นยำสูง สามารถสะท้อนความสนใจที่แท้จริงของผู้ใช้ และตอบโจทย์บริบทของธุรกิจเช่าอ่านหนังสือการ์ตูนในยุคดิจิทัลได้อย่างมีประสิทธิภาพ

% -----------------------------------------------------------
\section{วัตถุประสงค์}
% -----------------------------------------------------------

\subsection{เพื่อพัฒนาและสร้างเว็บแอปพลิเคชันต้นแบบ}
เพื่อพัฒนาและสร้างเว็บแอปพลิเคชันต้นแบบ ระบบแนะนำหนังสือการ์ตูน สำหรับเป็นส่วนติดต่อผู้ใช้งานและนำเสนอผลลัพธ์จากระบบแนะนำที่พัฒนาขึ้น

\subsection{เพื่อออกแบบและนำเสนอโมเดลระบบแนะนำ}
เพื่อออกแบบและนำเสนอโมเดลระบบแนะนำ ที่สามารถประมวลผลข้อมูลพฤติกรรมผู้ใช้ที่หลากหลาย เพื่อเพิ่มความแม่นยำในการแนะนำการเช่าอ่าน

\subsection{เพื่อประเมินและเปรียบเทียบประสิทธิภาพของโมเดล}
เพื่อประเมินและเปรียบเทียบประสิทธิภาพของโมเดลระบบแนะนำแบบหลายพฤติกรรมที่พัฒนาขึ้น กับระบบแนะนำแบบพฤติกรรมเดียว โดยใช้ตัวชี้วัดความแม่นยำ ได้แก่ Recall@K และ NDCG@K

% -----------------------------------------------------------
\section{ขอบเขตและข้อจำกัดของระบบ}
% -----------------------------------------------------------

\subsection{ขอบเขตด้านข้อมูล}

เนื่องจากระบบที่พัฒนาขึ้นใหม่ยังไม่มีประวัติการใช้งานจริง ในการทดสอบและเทรนโมเดล (Model Training) จึงใช้แนวทางดังนี้

\subsubsection{ชุดข้อมูล}
โมเดลหลักใช้ชุดข้อมูลมาตรฐานแบบ Multi-Behavior จากแหล่งข้อมูลสาธารณะโดยเว็บไซต์ MyAnimeList นำมาจำลองและจัดรูปแบบใหม่ให้สอดคล้องกับบริบทของการเช่าหนังสือการ์ตูน เพื่อใช้ฝึกโมเดลและวัดผลความแม่นยำ ส่วนโมเดลรอง ใช้ข้อมูลพฤติกรรมที่บันทึกในตาราง UserBehaviorLog ของเว็บแอปพลิเคชัน ได้แก่ การคลิกดูรายละเอียดหนังสือ การเพิ่มลงในตะกร้า และการเช่า เนื่องจากระบบยังไม่มีผู้ใช้งานจริงจำนวนมาก ข้อมูลที่ใช้ฝึกโมเดลรองในโครงงานนี้จึงเป็นข้อมูลจำลองของผู้ใช้ 1,000 คน ที่สร้างขึ้นตามการทำงานเดียวกับระบบจริง ซึ่งก็คือเพิ่มลงตะกร้า ยื่นคำขอเช่า อนุมัติ รับหนังสือ และคืนหนังสือ รวมกับผู้ใช้ทดสอบภายในระบบ รวมเป็นผู้ใช้ประมาณ 1,020 คน และหนังสือการ์ตูน 100 เรื่อง

\subsubsection{ตัวแปรพฤติกรรม}
ระบบจะนำข้อมูลพฤติกรรมของผู้ใช้มาวิเคราะห์ร่วมกันตามห่วงโซ่ 3 ขั้นตอน เพื่อค้นหาความสนใจแฝง ได้แก่ การคลิกอ่านเนื้อหา, การเพิ่มลงในตะกร้า และการทำรายการเช่า ซึ่งเป็นพฤติกรรมเป้าหมาย

\subsection{ขอบเขตด้านระบบงาน}

เป็นการพัฒนา Web Application ในลักษณะ Prototype เพื่อเป็นส่วนติดต่อผู้ใช้งานให้สามารถทดสอบระบบแนะนำได้จริง โดยแบ่งการทำงานเป็น 2 ส่วนหลัก คือ

\subsubsection{ส่วนหน้าเว็บแอปพลิเคชัน}
รองรับกระบวนการพื้นฐานของร้านเช่าหนังสือการ์ตูน ได้แก่ การสมัครสมาชิก การเลือกหนังสือการ์ตูนที่ชอบ 4 เรื่องหลังเข้าสู่ระบบครั้งแรกเพื่อแก้ปัญหา Cold-Start การค้นหาหนังสือ การดูรายละเอียดหนังสือ การเพิ่มลงในตะกร้า การยื่นคำขอเช่า โดยจองออนไลน์และชำระเงินที่ร้านเมื่อรับหนังสือ การยกเลิกคำขอ การดูประวัติการเช่า และการให้คะแนนหนังสือที่คืนแล้ว ผู้ใช้จะเห็นรายการหนังสือในส่วน "แนะนำสำหรับคุณ" ที่คำนวณจากโมเดลรุ่นที่ผู้ดูแลระบบเลือกใช้งาน พฤติกรรมการคลิกดูรายละเอียด การเพิ่มลงตะกร้า และการเช่า จะถูกบันทึกเพื่อใช้ฝึกโมเดลรุ่นถัดไป ผลการแนะนำจึงเปลี่ยนไปตามพฤติกรรมเมื่อมีการฝึกและเปิดใช้งานโมเดลใหม่ ไม่ได้เปลี่ยนทันทีในแต่ละคลิก ส่วนผู้ดูแลระบบสามารถจัดการหนังสือ สำเนาหนังสือ สมาชิก คำขอเช่า การคืนและค่าปรับ การเช่าหน้าร้าน ประวัติการเช่า การทดลองแบบ A/B Testing และการฝึกและสลับรุ่นโมเดล

\subsubsection{ส่วนประมวลผลระบบแนะนำ}
พัฒนาโมเดล Machine Learning แบบหลายพฤติกรรมด้วยเทคนิค MB-CGCN จำนวน 2 รุ่น ได้แก่ รุ่น v1 (Cart $\rightarrow$ Rent) ที่ฝึกจากชุดข้อมูล MyAnimeList และรุ่น v2 (Click $\rightarrow$ Cart $\rightarrow$ Rent) ที่ฝึกจาก Log พฤติกรรมภายในระบบ โมเดลถูกโหลดและประมวลผลภายในโปรเซสของ Backend เพื่อจัดอันดับหนังสือที่เหมาะสมมาแสดงผล

\subsection{ขอบเขตด้านเทคนิค}

\subsubsection{การพัฒนาโมเดล}
พัฒนาโมเดลด้วยภาษา Python โดยใช้ไลบรารีทางด้าน Deep Learning โดยใช้ PyTorch

\subsubsection{การพัฒนาเว็บแอปพลิเคชัน}
พัฒนาส่วน Backend ด้วย Django และ Django REST Framework และส่วน Frontend ด้วย React (Vite) ร่วมกับ Tailwind CSS ใช้ฐานข้อมูล PostgreSQL 15 เพียงระบบเดียวในการเก็บทั้งข้อมูลสมาชิก หนังสือ การเช่า และ Log พฤติกรรม ใช้ PyTorch Geometric ในการสร้างชั้น Graph Convolution และติดตั้งระบบด้วย Docker Compose สำหรับการพัฒนา และ Railway สำหรับการใช้งานจริง

\subsection{ข้อจำกัดของโครงงาน}

\subsubsection{ข้อจำกัดด้านข้อมูลเริ่มต้น}
แม้ระบบจะแก้ปัญหา Cold-Start ด้วยการให้ผู้ใช้เลือกหนังสือการ์ตูนที่ชอบ 4 เรื่องหลังเข้าสู่ระบบครั้งแรก เพื่อใช้สร้างโปรไฟล์ความสนใจเริ่มต้นและเริ่มแนะนำได้ทันที แต่ผลการแนะนำในระยะแรกอาจยังมีความแม่นยำไม่สูงนัก และจะดีขึ้นเรื่อย ๆ เมื่อระบบเก็บข้อมูลพฤติกรรมของผู้ใช้นั้น ๆ เพิ่มมากขึ้น

\subsubsection{ระบบการชำระเงิน}
ระบบไม่มีการชำระเงินออนไลน์ ผู้ใช้จองหนังสือผ่านเว็บและชำระค่าเช่าที่ร้านเมื่อมารับหนังสือ ระบบบันทึกเพียงยอดค่าเช่ารวม และค่าปรับ ของแต่ละคำสั่งเช่า และไม่มีการเชื่อมต่อกับธนาคารหรือระบบการเงินใด ๆ

\subsubsection{ข้อมูลสำหรับฝึกโมเดลรอง}
ข้อมูลพฤติกรรมที่ใช้ฝึกโมเดลรองส่วนใหญ่เป็นข้อมูลจำลอง และมีหนังสือการ์ตูนเพียง 100 เรื่อง ผลการประเมินโมเดลรองจึงสะท้อนรูปแบบของข้อมูลจำลอง และยังไม่ได้ยืนยันกับพฤติกรรมของผู้ใช้จริง

\subsubsection{ทรัพยากรการประมวลผล}
เนื่องจากการเทรนโมเดล Deep Learning ต้องใช้ทรัพยากรเครื่องสูง ขอบเขตของข้อมูลที่ใช้ทดสอบจึงอาจมีการจำกัดจำนวน ตามศักยภาพของฮาร์ดแวร์ที่ใช้ในการทำโครงงาน

% -----------------------------------------------------------
\section{ประโยชน์ที่คาดว่าจะได้รับ}
โครงงานนี้คาดว่าจะให้ประโยชน์ดังต่อไปนี้

\subsection{ประโยชน์สำหรับนักศึกษา}
ช่วยให้ศึกษาการพัฒนาระบบเว็บแบบองค์รวม การออกแบบและฝึกโมเดลระบบแนะนำแบบหลายพฤติกรรม และการประยุกต์ใช้ข้อมูลจริงในการแก้ปัญหา

\subsection{ประโยชน์สำหรับผู้ใช้ทั่วไป}
ช่วยให้ผู้ใช้ค้นพบหนังสือการ์ตูนที่ตรงกับความสนใจได้ง่ายขึ้น ลดเวลาในการค้นหา และได้รับการแนะนำที่เหมาะสมกับพฤติกรรมการใช้งาน

\subsection{ประโยชน์สำหรับเจ้าของแพลตฟอร์มหรือผู้ดูแลระบบ}
ช่วยเพิ่มอัตราการมีส่วนร่วมของผู้ใช้ สนับสนุนการตัดสินใจในการบริหารจัดการเนื้อหา และเก็บข้อมูลพฤติกรรมเพื่อนำไปพัฒนาระบบต่อไป

\subsection{ประโยชน์ทางวิชาการ}
เป็นกรณีศึกษาการประยุกต์ใช้โมเดล MB-CGCN กับข้อมูลพฤติกรรมหลายประเภทในโดเมนการเช่าอ่านและเป็นเอกสารอ้างอิงสำหรับงานวิจัยที่เกี่ยวข้อง

\subsection{ประโยชน์ทางเทคโนโลยี}
นำเสนอการผสานระหว่างระบบเว็บ Django REST Framework กับโมเดลแนะนำแบบกราฟที่โหลดไว้ในโปรเซสของเซิร์ฟเวอร์ ทำให้ตอบคำขอแนะนำได้ทันทีโดยไม่ต้องมีเซิร์ฟเวอร์ประมวลผลโมเดลแยก พร้อมกลไกฝึกโมเดลใหม่และสลับรุ่นโมเดลจากหน้าผู้ดูแลระบบ
